% ==========================================================
% CAPÍTULO 2 — TECNOLOGIAS UTILIZADAS
% ==========================================================
\chapter{Tecnologias Utilizadas}

Nesta seção são descritas as principais tecnologias e ferramentas empregadas até o momento no desenvolvimento do projeto.

\section{Front-end}
O front-end da aplicação está sendo construído com o \textbf{React 19}, utilizando \textbf{TypeScript} para garantir tipagem estática e maior segurança no desenvolvimento. O bundler \textbf{Vite} (com o plugin oficial \texttt{@tailwindcss/vite}) cuida do ambiente de desenvolvimento e da etapa de build. O estilo das telas é feito predominantemente com \textbf{Tailwind CSS 4}, configurado a partir do arquivo \texttt{frontend/src/styles/globals.css}.

\subsection{Principais bibliotecas}
\begin{itemize}
    \item \textbf{React Router DOM (v7)}: realiza o roteamento entre as páginas públicas (landing, login, cadastro) e áreas protegidas (home, perfil);
    \item \textbf{React Icons}: provê os ícones utilizados em formulários e cabeçalhos, garantindo consistência visual sem importar SVGs manualmente;
    \item \textbf{Fetch API}: encapsulada no utilitário \texttt{frontend/src/libs/api.ts}, centralizando headers, tratamento de erros e inclusão condicional do token JWT;
    \item \textbf{ESLint + TypeScript ESLint}: padronizam o estilo e a qualidade do código durante o desenvolvimento (\texttt{frontend/eslint.config.js}).
\end{itemize}

\section{Back-end}
O back-end está sendo construído com o framework \textbf{Flask 3}, escrito em Python, responsável pela definição das rotas e pela lógica de autenticação. O módulo \textbf{Flask-CORS} libera o domínio do front-end durante o desenvolvimento, enquanto a biblioteca \textbf{PyJWT} gera e valida os tokens JWT usados na autenticação do usuário. A configuração de variáveis sensíveis é feita com \textbf{python-dotenv}, e o acesso ao banco utiliza \textbf{mysql-connector-python} com pool de conexões (arquivo \texttt{backend/db.py}).

\section{Banco de Dados}
Adotamos o \textbf{MySQL} como Sistema Gerenciador de Banco de Dados, armazenando informações de autenticação na tabela \texttt{Usuario}. O acesso é realizado por meio de um pool de conexões configurado em \texttt{backend/db.py}, reduzindo a sobrecarga na abertura de conexões a cada requisição. A camada de serviço em \texttt{backend/app.py} executa consultas parametrizadas para evitar injeção de SQL e facilitar a evolução futura do schema.

\section{Controle de Versão}
O versionamento do projeto está sendo realizado através do \textbf{Git} e hospedado no \textbf{GitHub}, permitindo o acompanhamento de alterações, controle de branches e colaboração entre os membros do grupo.

\section{IDE e Ferramentas de Desenvolvimento}
A principal IDE utilizada é o \textbf{Visual Studio Code}, configurado com extensões para suporte ao React, TypeScript, formatação de código e integração direta com o GitHub. O \textbf{Node.js} (via \textit{npm}) gerencia as dependências do front-end, enquanto o ambiente Python é gerenciado por um \texttt{venv} descrito no arquivo \texttt{backend/requirements.txt}.
